% UBT-AI-PROVENANCE-BEGIN
% schema: ubt-ai-provenance/v1
% tier: C_working
% ai_assistance: disclosed
% human_review: risk-based
% editorial_responsibility: Ing. David Jaroš
% policy: ../../../AI_PROVENANCE.md
% notice: Working material; exhaustive human review is not claimed.
% UBT-AI-PROVENANCE-END

\chapter{Lorentzova geometrie na hermitovském řezu \(\mathbb H_{\mathbb C}\)}
\label{ch:lorentz-hc}

\section{Klasický základ: \(\mathrm{Herm}(2,\mathbb C)\)}
Tato sekce je standardní spinor/Lorentz matematika.  Je součástí, protože
UBT používá stejnou algebru prostřednictvím izomorfismu
\[
 \mathbb B=\mathbb C\otimes\mathbb H\cong M_2(\mathbb C).
\]
Nechť $\sigma_i$ jsou Pauliho matice a spojují skutečný čtyřvektor
$x=(x^0,x^1,x^2,x^3)$ s hermitovskou maticí
\[
 X=x^0 I_2+x^i\sigma_i
 =\begin{pmatrix}
 x^0+x^3 & x^1-i x^2\\
 x^1+i x^2 & x^0-x^3
 \end{pmatrix}.
\]
Výpočet přímé determinanty dává
\[
 \boxed{\det X=(x^0)^2-(x^1)^2-(x^2)^2-(x^3)^2.}
\]
Takže determinant na hermitovském řezu je přesně Minkowského kvadratika
formulář s podpisem $(+---)$.

\section{Akce \(SL(2,\mathbb C)\)}
Pro $A\in SL(2,\mathbb C)$ definujte
\[
 X\longmapsto X'=A X A^\dagger .
\]
Poustevnost je zachována, protože
\[
 (A X A^\dagger)^\dagger=A X A^\dagger.
\]
Zachován je také determinant:
\[
 \det X'=\det A\,\det X\,\det A^\dagger
 =\det X,
\]
od $\det A=1$ a $\det A^\dagger=\overline{\det A}=1$.
Proto indukovaná reálná lineární transformace čtyř koeficientů
zachovává Minkowského normu.  Nepřetržité omezování na identitu dává a
homomorfismus
\[
 SL(2,\mathbb C)\longrightarrow SO^+(1,3).
\]
$A$ i $-A$ působí identicky na $X$, takže jádro obsahuje
$\{\pm I_2\}$.  Standardním výsledkem je dvojitý kryt
\[
 \boxed{SL(2,\mathbb C)/\{\pm I_2\}\cong SO^+(1,3).}
\]
Pro tuto konstrukci není vyžadována žádná algebra dělené čtveřice.

\section{Převod do bikvaternionové notace}
Použijte kanonickou maticovou reprezentaci
\[
 \mathbf e_i\longmapsto i\sigma_i.
\]
Pak $\sigma_i\leftrightarrow -i\mathbf e_i$ a hermitovská matice výše
odpovídá biquaternionu
\[
 X_{\mathbb B}=x^0\mathbf1-i x^i\mathbf e_i.
\]
Toto je prostě další popis souřadnic téhož
Plátek $\mathrm{Herm}(2,\mathbb C)$.

Současná tetradská konvence UBT místo toho často používá Lorentzův reálný základ
\[
 \mathbf u_0=i\mathbf1,
 \qquad \mathbf u_k=\mathbf e_k,
\]
tak je noha tetrády
\[
 E_\mu=i e_\mu{}^0\mathbf1+e_\mu{}^k\mathbf e_k.
\]
S kvaternionovou konjugací $\sharp$,
\[
 E_\mu^\sharp=i e_\mu{}^0\mathbf1-e_\mu{}^k\mathbf e_k.
\]
Pomocí $\mathbf e_j\mathbf e_k+\mathbf e_k\mathbf e_j=-2\delta_{jk}$,
člověk získá
\begin{align*}
 \frac12(E_\mu^\sharp E_\nu+E_\nu^\sharp E_\mu)
 &=\left(-e_\mu{}^0e_\nu{}^0+e_\mu{}^k e_\nu{}^k\right)\mathbf1\\
 &=g_{\mu\nu}\mathbf1.
\end{align*}
Toto je stejná Lorentzova geometrie s opačným celkovým podpisem
konvence $(-+++)$.  Podpisová konvence je vedení účetnictví; není to novinka
fyzický stupeň svobody.

\section{Proč je to důležité pro UBT}
Klasický výsledek stanoví tři užitečná fakta před jakoukoli dynamikou UBT
představil:
\begin{enumerate}
 \item Komplexní kvaterniony již nesou spinorální Lorentzovu akci;
 \item ze čtyřrozměrné Lorentzovy kvadratické formy lze číst interně
       algebra;
 \item levé/pravé násobení záleží, protože Lorentzova akce ano
       přirozeně oboustranné v maticovém jazyce.
\end{enumerate}

To, co tato klasická konstrukce\emph{není} dokazuje, je stejně důležité.
Nedokazuje to, že hlavní pole UBT dynamicky generuje každé přání
tetráda, neodvozuje Einsteinovy rovnice a nevybírá
fyzické spojení.  Toto jsou problémy s mosty specifické pro UBT, které jsou řešeny v
kovariantní tetrádová kapitola a kanonická závěrečná kniha GR.

\section{Konkrétní boost jako kontrolní příklad}
Zvyšte rychlost $\chi$ podél třetí prostorové osy,
\[
 A=\exp\left(\frac{\chi}{2}\sigma_3\right)
 =\begin{pmatrix}e^{\chi/2}&0\\0&e^{-\chi/2}\end{pmatrix}.
\]
Použití $X'=AXA^\dagger$ a porovnání koeficientů dává
\[
 \begin{pmatrix}x'^0\\x'^3\end{pmatrix}
 =\begin{pmatrix}\cosh\chi&\sinh\chi\\
                  \sinh\chi&\cosh\chi\end{pmatrix}
  \begin{pmatrix}x^0\\x^3\end{pmatrix},
 \qquad x'^1=x^1,\quad x'^2=x^2.
\]
Pak
\[
 (x'^0)^2-(x'^3)^2=(x^0)^2-(x^3)^2,
\]
což je obyčejná Lorentzova podpora napsaná jako akce vnitřní algebry.

\section{Živé zdroje a historická poznámka}
Konvence algebry je zachována
\texttt{canonical/algebra/biquaternion\_algebra.tex}; současné použití
Lorentz-real tetrad je zdokumentován v aktivní GR stopě.  Bývalý
\texttt{appendix\_P6\_lorentz\_in\_HC.tex} v GR zůstává v historickém archivu,
ale živá učebnice ji již neimportuje.
